\documentclass[10pt,a4paper]{amsart}
\usepackage[margin=19mm]{geometry}
\usepackage{amsmath,amssymb,mathtools}
\usepackage[scr=rsfs,cal=euler]{mathalfa}
\usepackage{xfrac}
\usepackage[unicode,hidelinks,bookmarks=false]{hyperref}
\newcommand{\ie}{i.e.\ }
\newcommand{\cf}{cf.\ }
\newcommand{\resp}{respectively\ }
\newcommand{\Subsection}{\subsection}
\providecommand{\nobreakdash}{\nobreak\hbox{-}}
\setlength{\emergencystretch}{2em}
\multlinegap0pt
\usepackage{amssymb}
\usepackage[all]{xy}
\usepackage{mathtools}
\usepackage{booktabs}
\usepackage{float}
\usepackage{caption}
\newtheorem{corollary}{Corollary}
\newtheorem{lemma}{Lemma}
\newtheorem{proposition}{Proposition}
\newcommand{\eqst}[1]{\begin{equation*}#1\end{equation*}}
\title{Arithmetic aspects of discrete periodic Toda flows}
\author{Bora Yalkinoglu}
\date{Source: arXiv:2408.07315v3 (CC BY 4.0)}
\begin{document}
\maketitle

\noindent\emph{Source and licence.} Arithmetic aspects of discrete periodic Toda flows. Version arXiv:2408.07315v3 (CC BY 4.0), distributed by arXiv under the Creative Commons Attribution 4.0 International licence, http://creativecommons.org/licenses/by/4.0/, as recorded in the arXiv licence metadata for this exact version and shown on the abstract page https://arxiv.org/abs/2408.07315v3. Full licence notice: \texttt{licenses/arxiv-2408.07315-CC-BY-4.0.txt}.\par\smallskip

\noindent\textit{Editorial scope.} Counted unit: the article's proposition cyclictime (source0.tex lines 624--632). The article's complete original proof of it is delivered with it (source0.tex lines 633--647, one \texttt{proof} environment of 15 lines). No mathematics is written, repaired or re-derived: the statement and the proof are the article's own lines, in the article's own notation, and no retained line is substituted or re-flowed.  Retained uncounted as the source-local context the proof needs: lines 416--419; lines 450--454.  Nothing was omitted from any retained span, so the package carries no omission record.  No retained line contains a citation, so the wrapper carries no bibliography and no bibliography entry.  No retained line contains a figure environment, an image-inclusion command or drawing code, so no image is delivered and the package declares no asset dependency.  Editorial formatting only, in addition: an AMS \texttt{amsart} preamble that restates the article's own declarations and macros used by the retained lines, namely the environments \texttt{corollary}, \texttt{lemma}, \texttt{proposition} on separate counters, together with the standard packages those lines need.  The retained environments print in this document as Corollary 1, Lemma 1, Proposition 1 in order of appearance, whereas the article prints the same items with the numbering of its own full document.  The complete native source of the article as distributed in the arXiv e-print for this exact version is bundled unchanged as \texttt{sources/163765/source0.tex}.
\begin{corollary}
\label{easycor}
We have \eqst{w^\sigma = v.}
\end{corollary}

\begin{lemma}
\label{keyrelation}
The following relations hold
\begin{align*}I_n u - v &= I_n u_{V_{n}} \\ &= I_n \overline u - \overline w, \\ (V_{n-1}+x)u-w &= (V_n + x)\overline u - \overline v.\end{align*}
\end{lemma}

\begin{proposition}
\label{cyclictime}
The following relations hold
\begin{align*}
x \overline u &= x u +V_{n-1} u_{I_{n-1}} + v - V_n u^\sigma \\
			&= h +2v - I_n u  - V_n u^\sigma, \\
\overline v &= V_n \overline u + v - V_n u^\sigma.
\end{align*}
\end{proposition}

\begin{proof}
For the first assertion we calculate 
\begin{align*} x \overline u &= (\tilde v ^2 + \tilde v h +f)/(-I_n u) \\
&= (\tilde v ((V_{n-1}+x) u - w) + f) / (-I_n u) \\
&= ((-I_n u + v) (V_{n-1}+x) u - (-I_n u +v) w + f)/(-I_n u) \\
&= (V_{n-1}+x) u - w + ((V_{n-1}+x) u v - vw + f)/(-I_n u).
\end{align*}
On the other hand we have \begin{align*}
V_n u^\sigma &= (v^2 + vh +f)/(I_nu) \\
&=  ((I_n+V_{n-1}+x)uv-vw +f) / (I_n u) \\
&= v + ((V_{n-1}+x)uv-vw+f)/(I_n u).
\end{align*}
The result now follows with the observation $V_{n-1}u_{I_{n-1}} = V_{n-1}u-w$.
The second assertion follows from Lemma \ref{keyrelation} and Corollary \ref{easycor}.
\end{proof}

\end{document}
